\section{The Architect's Question}\label{the-architects-question}

After this chapter we should be able to take a vague deployment ask ---
``thousands of employees asking questions over internal documents'' ---
and turn it into a \emph{measurable workload} across six concrete
dimensions. We will trace the arithmetic from user concurrency through
traffic, token profile, latency, economics, and operational constraints,
ending with a workload characterization that directly enables system
architecture decisions. After this chapter, the question ``which model
should we choose?'' can be answered with numbers, not intuition.

\section{Concept}\label{concept}

\subsection{The Six-Dimension Workload-Characterization
Framework}\label{the-six-dimension-workload-characterization-framework}

A workload is defined by six orthogonal dimensions. Together they form a
complete specification that --- taken alone --- determines model
selection, system architecture, and economic feasibility. No single
dimension is sufficient; the architect must characterize all six before
committing to a design.

\begin{enumerate}
\def\labelenumi{\arabic{enumi}.}
\item
  \textbf{Quality} --- the fidelity, accuracy, and capability the
  workload demands of the model. Includes reasoning depth, tool-use
  requirements, modality (text‑only, multimodal, etc.), and any accuracy
  thresholds (e.g.~\textgreater85\% factual recall on internal Q\&A).
  Quality is the \emph{goal}; the other dimensions are the
  \emph{constraints} within which quality must be delivered.
\item
  \textbf{Traffic} --- the request rate the system must sustain,
  expressed in requests per second (rps) or queries per second, plus
  concurrency levels and burstiness. Traffic determines the compute and
  memory bandwidth required to keep the serving pipeline saturated
  without excessive queuing.
\item
  \textbf{Token profile} --- the distribution of input length, output
  length, and context length per request. This is the workload's
  ``shape'' in token space: how many tokens arrive in the prompt, how
  many the model generates, and whether the prompt is dense (long
  context) or sparse. The token profile is the primary driver of memory
  (KV cache) and prefill cost.
\item
  \textbf{Latency} --- the time budget for different stages of request
  processing. Typically split into \emph{time-to-first-token} (TTFT),
  which encompasses retrieval and prefill, and
  \emph{time-per-output-token} (TPOT), which governs decode length.
  Latency budgets feed directly into SLOs and serve as the key
  differentiator between serving configurations (e.g.~continuous
  batching vs.~discrete batching).
\item
  \textbf{Economic constraints} --- the cost budget denominated in
  tokens per dollar, dollars per million tokens, or cost per request.
  This dimension translates traffic and token profile into a recurring
  cost line, and it determines whether a given infrastructure choice
  (single node vs.~multi‑node, FP16 vs.~FP8) is affordable at scale.
\item
  \textbf{Operational constraints} --- availability, privacy, locality,
  and update frequency. Multi‑region deployment for fault tolerance.
  Data‑privacy requirements that force on‑prem or VPC‑local embedding
  models. Update frequency that dictates how often embeddings or model
  weights must be refreshed. These constraints often interact with
  traffic and economics to force architectural compromises.
\end{enumerate}

The framework is deliberately dimensional: a workload that is
``high‑quality, high‑traffic, long‑context, low‑latency, cost‑sensitive,
private'' is fundamentally different from ``low‑quality, low‑traffic,
short‑context, high‑latency, cost‑abundant, public,'' even if the raw
text looks the same. The framework prevents the architect from
optimizing one dimension at the expense of others --- a common failure
mode.

\section{Mental Model}\label{mental-model}

Think of the six dimensions as the axes of a six‑dimensional space in
which every AI workload resides. A workload's position in this space is
its \emph{fingerprint}. When an architect says ``we need a model for
enterprise Q\&A,'' that is only a descriptor of the quality dimension.
To make the statement actionable, we must project the workload onto all
six axes: \emph{what quality level, how much traffic, what token
profile, what latency SLO, what economic ceiling, what operational
constraints?} The intersection of these projections is the workload
point that drives every downstream decision --- model selection,
infrastructure sizing, serving configuration, and TCO.

The mental model is not a checklist; it is a \emph{lens}. Looking
through it, the architect sees which dimensions are tight (binding) and
which are loose (permissive). Tight dimensions become the design
drivers; loose dimensions offer optimization freedom. The goal is to
identify the binding constraints so that resources are allocated where
they matter most.

\section{Worked Example}\label{worked-example}

\subsection{Characterizing the Canonical Enterprise-Q\&A RAG
Workload}\label{characterizing-the-canonical-enterprise-qa-rag-workload}

We now apply the six‑dimension framework to the \textbf{canonical
enterprise‑Q\&A RAG workload} defined in §14 (book-architecture.md). The
canonical numbers are the fixed reference set for Part II; all
arithmetic in this chapter traces to them.

\subsubsection{\hyperref[tab:4.1]{Table 4-1} --- Six-dimension workload characterization for
the canonical RAG
workload}\label{table-4-1-six-dimension-workload-characterization-for-the-canonical-rag-workload}

\begin{longtable}[]{@{}
  >{\raggedright\arraybackslash}p{0.0833\linewidth}
  >{\raggedright\arraybackslash}p{0.3125\linewidth}
  >{\raggedright\arraybackslash}p{0.6042\linewidth}@{}}
\toprule\noalign{}
\begin{minipage}[b]{\linewidth}\raggedright
Dimension
\end{minipage} & \begin{minipage}[b]{\linewidth}\raggedright
Characteristic
\end{minipage} & \begin{minipage}[b]{\linewidth}\raggedright
Derivation / Source
\end{minipage} \\
\midrule\noalign{}
\endhead
\bottomrule\noalign{}
\endlastfoot
\textbf{Quality} & Factual accuracy \textgreater85\% on enterprise Q\&A;
reasoning via retrieval; no tool use beyond search API & {[}2°{]}
DERIVED from enterprise benchmark suites (e.g.~HotpotQA‑style retrieval
Q\&A) \\
\textbf{Traffic} & \textasciitilde10 rps average; peaks
\textasciitilde40 rps & \textbf{Derivation}: 2,000 registered users ×
5\% concurrency = 100 concurrent users. By Little's Law (L = λW), with L
= 100 and average request duration W ≈ 10 s, throughput λ = L/W = 100/10
≈ 10 rps. Peaks \textasciitilde40 rps arise when concurrency spikes to
\textasciitilde20\% (400 users) with the same 10 s duration, yielding λ
= 400/10 = 40 rps. \\
\textbf{Token profile} & Input: \textasciitilde1,200 prompt +
\textasciitilde8,000 retrieved context ≈ 9,200 tokens; Output:
\textasciitilde300 tokens; total ≈ 9,500 tokens/request & {[}1P{]} §14
canonical scenario; tokenizer‑verified on the target model's
tokenizer \\
\textbf{Latency} & TTFT budget 1.2 s (retrieval \textasciitilde120 ms +
prefill \textasciitilde1.08 s); TPOT budget \textasciitilde25 ms/token;
p95 TTFT ≤ 2 s, p95 TPOT ≤ 35 ms & {[}2°{]} SLO-derived from
user‑experience targets; retrieval latency from vector DB on same‑region
deployment \\
\textbf{Economic constraints} & \textasciitilde\$1.20 per 1M input
tokens, \textasciitilde\$2.00 per 1M output tokens on 8×H100 cloud
instance; \textasciitilde95,000 input tokens/s per dollar at 10 rps;
infrastructure cost ≈ \$3.50/hour for the host & {[}VERIFY{]} DERIVED
from cloud GPU pricing as of 2026 (on‑demand H100 instances); tokens/s
per dollar = total token throughput ÷ per-second cost (\$3.50/hr ÷
3600) \\
\textbf{Operational constraints} & Multi‑region deployment
(active‑active for availability); embeddings refreshed daily from
document store; privacy‑sensitive documents force on‑prem or VPC‑local
retrieval; 99.9\% availability SLA & {[}2°{]} DERIVED from typical
enterprise IT policy and the canonical scenario's availability
requirements \\

\label{tab:4.1}\end{longtable}

\subsubsection{\hyperref[tab:4.2]{Table 4-2} --- Canonical workload characterization across
six
dimensions}\label{table-4-2-canonical-workload-characterization-across-six-dimensions}

\begin{longtable}[]{@{}
  >{\raggedright\arraybackslash}p{0.3333\linewidth}
  >{\raggedright\arraybackslash}p{0.3333\linewidth}
  >{\raggedright\arraybackslash}p{0.3333\linewidth}@{}}
\toprule\noalign{}
\begin{minipage}[b]{\linewidth}\raggedright
Metric
\end{minipage} & \begin{minipage}[b]{\linewidth}\raggedright
Value
\end{minipage} & \begin{minipage}[b]{\linewidth}\raggedright
Derivation / Source
\end{minipage} \\
\midrule\noalign{}
\endhead
\bottomrule\noalign{}
\endlastfoot
Quality & \textgreater85\% factual accuracy {[}2° DERIVED{]} &
Enterprise Q\&A benchmarks \\
Traffic & 10 rps avg; 40 rps peak {[}2° DERIVED{]} & 2,000 users × 5\%
concurrency = 100 concurrent; 100/10s = 10 rps; peaks at 20\%
concurrency → 40 rps \\
Token profile & 9,200 input tokens + 300 output tokens {[}1P §14{]} &
1,200 prompt + 8K context; 300‑token answer \\
Latency & TTFT 1.2 s (retrieval \textasciitilde120 ms + prefill); TPOT
\textasciitilde25 ms/token {[}2° SLO{]} & User‑experience targets \\
Economics & \$1.20/M input; \$2.00/M output {[}VERIFY DERIVED{]} & 2026
on‑demand H100 pricing; tokens/s per dollar \\
Operational & Multi‑region; daily embedding refresh; 99.9\% availability
{[}2° DERIVED{]} & Enterprise IT policy \\

\label{tab:4.2}\end{longtable}

\emph{(All figures trace to the §14 canonical scenario; none are
independent measurement claims. \hyperref[tab:4.2]{Table 4-2} consolidates the six‑dimension
characterization for quick reference.)}

\subsubsection{Arithmetic Walk‑Through: From Users to Requests per
Second}\label{arithmetic-walkthrough-from-users-to-requests-per-second}

The step from ``\textasciitilde2,000 registered users'' to
``\textasciitilde10 rps average'' is the key connective tissue that
makes the continuous scenario concrete. Here is the full derivation,
traceable line by line:

\begin{enumerate}
\def\labelenumi{\arabic{enumi}.}
\tightlist
\item
  \textbf{Concurrency}: a fraction \(f\) of the \(N\) registered users
  are actively in flight at any moment. With \(f = 0.05\) and
  \(N = 2{,}000\):
\end{enumerate}

\[
C = f \cdot N = 0.05 \times 2{,}000 = 100 \text{ concurrent users}
\]

\begin{enumerate}
\def\labelenumi{\arabic{enumi}.}
\setcounter{enumi}{1}
\item
  \textbf{Request completion time} (service time \(W\)): the average
  end‑to‑end time from request submission to answer delivery is
  \textasciitilde10 s. This comprises retrieval (\textasciitilde120 ms)
  + prefill (\textasciitilde1.08 s for 9,200 tokens on H100) + decode of
  300 tokens at \textasciitilde25 ms/token (\textasciitilde7.5 s). The
  dominant term is decode, but the full pipeline sits at
  \textasciitilde10 s.
\item
  \textbf{Throughput via Little's Law}: Little's Law, \(L = \lambda W\),
  relates average concurrency \(L\), average arrival rate \(\lambda\),
  and average time in system \(W\). Rearranging to solve for throughput:
\end{enumerate}

\[
\lambda = \frac{L}{W} = \frac{100 \text{ concurrent}}{10 \text{ s}} = 10 \text{ rps}
\]

\begin{enumerate}
\def\labelenumi{\arabic{enumi}.}
\setcounter{enumi}{3}
\tightlist
\item
  \textbf{Peak throughput}: if concurrency spikes to \(f = 0.20\)
  (i.e.~400 users) and the per‑request duration stays \(W = 10\) s,
  Little's Law gives
\end{enumerate}

\[
\lambda_{\text{peak}} = \frac{f_{\text{peak}} \cdot N}{W} = \frac{0.20 \times 2{,}000}{10} = 40 \text{ rps}
\]

This matches the canonical peak of \textasciitilde40 rps.

The derivation is explicit: the ``\textasciitilde5\% concurrently
active'' and ``\textasciitilde10 rps average'' are not independently
asserted; they are linked by the 10 s average request duration. Change
any one number and the others shift accordingly. This is the purpose of
the exercise --- not to produce a fixed set of gospel numbers, but to
establish \emph{how} the numbers connect, so the architect can recompute
them when the requirement changes.

Before moving on, it is worth naming the four distinct quantities so
they are not collapsed: \textbf{registered users} are the population;
\textbf{concurrent users} are the activity state (those in flight at a
moment); \textbf{arrival rate} (λ) is requests per unit time; and
\textbf{service time} (W) is the end-to-end duration of a request. They
are related by Little's Law, C = λW, but the architect must
\emph{estimate or measure} concurrency, arrival rate, and service time
--- none is a fixed property of the workload. The ``5\% active'' and
``10 s duration'' are deliberately simple scenario assumptions
({[}ILLUSTRATIVE{]}), not externally verified facts; in a real
engagement the architect replaces them with observed concurrency,
measured latency, and a burst profile from telemetry. Getting this
separation right is what turns traffic modeling from a guess into a
reproducible arithmetic.

\subsubsection{Token Profile in Detail}\label{token-profile-in-detail}

The canonical workload's token profile is input‑heavy, which has direct
consequences for system design:

\begin{itemize}
\tightlist
\item
  \textbf{Per‑request input}: 1,200 tokens (enterprise query, possibly
  reformulated) + 8,000 tokens (retrieved context from vector DB) ≈
  9,200 tokens. The {[}1P{]} provenance traces to the §14 canonical
  scenario.
\item
  \textbf{Per‑request output}: \textasciitilde300 tokens (the generated
  answer, possibly with citations).
\item
  \textbf{Input‑to‑output ratio}: 9,200 ÷ 300 ≈ 30× more input tokens
  than output tokens. This ratio is {[}2° DERIVED{]} from the canonical
  numbers and is the single biggest factor in why this workload is
  memory‑bound (KV cache) rather than decode‑bound.
\item
  \textbf{At 10 rps}: input tokens/s = 9,200 × 10 = 92,000 tokens/s;
  output tokens/s = 300 × 10 = 3,000 tokens/s. At peak 40 rps, input
  spikes to \textasciitilde368,000 tokens/s and output to
  \textasciitilde12,000 tokens/s. These {[}2° DERIVED{]} numbers appear
  in \hyperref[tab:4.2]{Table 4-2}.
\end{itemize}

The input‑heavy profile means that \emph{prefill} (processing the
prompt) dominates the latency and cost budget. A model that processes 8K
tokens of context in under 1 s of prefill is essential; otherwise the
TTFT budget of 1.2 s cannot be met. This is why the token profile is the
primary architectural driver for this workload.

\subsubsection{Latency from SLO}\label{latency-from-slo}

The latency dimension is specified through Service Level Objectives,
which translate user‑experience goals into concrete time budgets:

\begin{itemize}
\tightlist
\item
  \textbf{TTFT (time-to-first-token)}: budget of 1.2 s, comprising
  retrieval (\textasciitilde120 ms for vector search and reranking on a
  local GPU) + prefill (\textasciitilde1.08 s for 9,200 tokens on an
  H100, assuming \textasciitilde8.5K tokens/s prefill throughput). p95
  TTFT must stay ≤ 2 s, allowing headroom for occasional cache misses or
  network jitter.
\item
  \textbf{TPOT (time-per-output-token)}: budget of \textasciitilde25
  ms/token. A 300‑token answer therefore takes \textasciitilde7.5 s of
  total decode time. p95 TPOT ≤ 35 ms/token provides headroom for
  batching variability.
\end{itemize}

These budgets are {[}2° DERIVED{]} from typical enterprise Q\&A user
expectations (sub‑2‑second feel) and from the hardware's published
prefill/decode rates. They are the latency constraints that every
subsequent architectural decision must respect.

\subsubsection{Economic Constraints}\label{economic-constraints}

The economic dimension translates the token profile and traffic into a
cost structure:

\begin{itemize}
\tightlist
\item
  \textbf{Infrastructure}: A single host with 8 × H100 GPUs (640 GB
  total GPU memory, 140 GB model FP16 weights fit with room for KV
  cache). Capital cost ≈ \$3.50/hour on-demand, or
  \textasciitilde\$2,500/month reserved.
\item
  \textbf{Token pricing}: \textasciitilde\$1.20 per million input
  tokens, \textasciitilde\$2.00 per million output tokens on the same
  H100 instance (derived from cloud provider pricing as of 2026).
\item
  \textbf{Throughput per dollar}: At 10 rps average, the system
  processes \textasciitilde92,000 input tokens/s + \textasciitilde3,000
  output tokens/s. Dividing by the \$3.50/hour infrastructure cost (≈
  \$0.00097 per second) yields \textasciitilde95,000 input tokens/s per
  dollar and \textasciitilde3,100 output tokens/s per dollar. These
  {[}VERIFY DERIVED{]} numbers are the economics framing the canonical
  scenario. (\hyperref[chap:5]{Chapter 5} expresses the same economics on a GPU list-price
  basis as \textbf{tokens per dollar-hour}; see its unit-reconciliation
  note before cross-chapter comparison.)
\item
  \textbf{Cost per request}: At 10 rps, each request carries
  \textasciitilde9,500 tokens (9,200 input + 300 output). At the
  per‑million rates, cost per request ≈ (\$1.20 × 9.2 + \$2.00 × 0.3) /
  1,000 ≈ \$0.015 per request per inference cycle. At 40 rps peak, cost
  scales linearly.
\end{itemize}

The economic constraint is what makes the workload real: a 70B FP16
model on one host can serve the canonical workload at the target SLO,
but scaling to higher traffic or longer contexts would require
additional hosts, and the cost line must be re‑evaluated.

\section{Measurement}\label{measurement}

For this chapter, measurement is about \textbf{quantifying the six
dimensions} so the workload can be communicated definitively and used to
drive architecture decisions. Three practical habits anchor the
architect:

\begin{enumerate}
\def\labelenumi{\arabic{enumi}.}
\item
  \textbf{Measure tokens, not words.} Run the model's own tokenizer on
  representative prompts from real traffic. The ``4 chars ≈ 1 token''
  heuristic is for estimation only; real counts differ by language,
  formatting, code, and tokenizer version. Input token counts directly
  determine KV cache size and prefill time; output token counts
  determine decode bandwidth.
\item
  \textbf{Split input and output.} Measure both legs of the request
  separately (prompt tokens and generated tokens), because they land on
  different bottlenecks --- input on memory/prefill, output on decode
  --- and on different cost line items. An input‑heavy profile (like
  RAG's 30× ratio) means prefill dominates; an output‑heavy profile
  would dominate decode.
\item
  \textbf{Log the distribution, not just the mean.} A workload that
  averages 9.2 K input tokens but has a long tail (e.g.~32 K peak
  contexts) has a very different KV cache and latency profile than one
  with a tight distribution. Log percentiles (p90, p99) alongside the
  average.
\end{enumerate}

These measurement habits are the token-layer answer to the book's
recurring question, ``what would I actually measure here?'' --- we
measure token counts and their distribution, at the edge, before any
architecture decision is made.

\section{Common Mistakes}\label{common-mistakes}

\begin{itemize}
\tightlist
\item
  \textbf{Treating ``5\% concurrency'' as the throughput}. Concurrency
  is a snapshot of simultaneous users; throughput depends on how long
  each request takes. 100 concurrent users with 200 s per request yield
  0.5 rps, not 10 rps. The binding connection is through request
  duration (Little's Law). Catch this by always asking ``how long is
  each request in system?'' before quoting a concurrency-derived
  throughput.
\item
  \textbf{Assuming ``one token ≈ one word''}. Enterprise Q\&A prompts
  with code snippets, JSON payloads, or technical terminology can run
  2--5× the naive token estimate, silently inflating KV cache and cost.
\item
  \textbf{Ignoring the input--output asymmetry}. An input-heavy workload
  like RAG is memory-bound (prefill/KV cache), while an output-heavy
  workload is decode-bound. Optimizing decode throughput on an
  input-heavy workload moves the needle negligibly; the real win is
  prefill reduction or KV cache reuse.
\item
  \textbf{Quoting context window as free capacity}. The 128 K token
  window is an upper bound; using 9.2 K of it still costs for the full
  9.2 K in memory and prefill time. Keeping context shorter than
  necessary is not ``free.''
\item
  \textbf{Overlooking operational constraints}. Multi--region deployment
  adds replication cost; privacy requirements may force a more expensive
  on-prem embedding model. These are often the true binding constraints,
  not the per-token price.
\end{itemize}

\section{Architecture Consequence}\label{architecture-consequence}

The six‑dimension characterization directly dictates the architectural
path for the canonical enterprise‑Q\&A RAG workload:

\begin{figure}
\centering
\pandocbounded{\includegraphics[keepaspectratio,alt={\hyperref[fig:4.1]{Fig 4.1} --- Six-dimension workload characterization mapped to architectural decisions. {[}ILLUSTRATIVE conceptual{]}},width=\maxwidth{\textwidth},keepaspectratio]{figures/fig-04-0401.pdf}}
\caption{\hyperref[fig:4.1]{Fig 4.1} --- Six-dimension workload characterization mapped to
architectural decisions. {[}ILLUSTRATIVE conceptual{]}}

\label{fig:4.1}\end{figure}

\emph{\hyperref[fig:4.1]{Fig 4.1} --- The six workload dimensions and the architectural
decision each one drives: Quality → model size/type; Traffic →
concurrency \& batching strategy; Token profile → KV cache size \&
prefill demand; Latency → TTFT/TPOT targets \& batch window; Economic →
host count \& cost ceiling; Operational → multi-region vs.~single-region
deployment.}

\begin{itemize}
\tightlist
\item
  \textbf{Model selection}: A 70B FP16 dense model (140 GB weights) fits
  on a single host with 8 × H100 (640 GB GPU memory). The model is large
  enough to answer factual enterprise questions without fine‑tuning, but
  the 140 GB footprint means KV cache for 9.2 K context adds
  \textasciitilde20--30 GB of GPU memory per request at peak, leaving
  headroom but not abundance.
\item
  \textbf{Serving configuration}: One host is the baseline. Continuous
  batching (e.g.~vLLM) is nearly mandatory to achieve the 1.2 s TTFT
  budget under 10 rps input‑heavy traffic; without it, prefill of 9.2 K
  tokens per request would serialize and push TTFT well above 2 s. The
  input‑heavy token profile (30× more input than output) makes
  continuous batching especially effective, as many requests share the
  same prefix from retrieved context.
\item
  \textbf{Memory planning}: KV cache for 9.2 K context on the 70B FP16
  canonical model is 2 × layers × hidden × bytes ≈ 2.5 MB/token (\hyperref[chap:7]{Chapter 7}), giving ≈9,200 × 2.5 MB ≈ 23 GB per request. At 10 concurrent
  full-context requests that is \textasciitilde230 GB of KV on top of
  the 140 GB weights --- pressing against the 640 GB pool well before
  concurrency reaches 100. The architecture must therefore limit
  concurrency, quantize KV (FP8 → \textasciitilde1.3 MB/token,
  \textasciitilde12 GB/request halves it), or reduce context. This is
  the same KV-residency discipline developed fully in Chapters 7 and 17.
\item
  \textbf{Economic feasibility}: At \textasciitilde\$3.50/hour per host
  and \textasciitilde\$0.015 per request, the TCO is driven by the 9.2 K
  input tokens per request. If the input token count could be reduced to
  2 K (e.g.~via better retrieval or query expansion), cost per request
  drops to \textasciitilde\$0.004, and the same traffic fits within a
  much lower budget. This is the lever the architect pulls when TCO is
  the binding constraint.
\item
  \textbf{Operational topology}: Multi‑region deployment (active‑active)
  provides availability but doubles the infrastructure cost. If the
  99.9\% availability SLA is non‑negotiable, the architecture must
  absorb the 2× cost. If it is negotiable, a single-region with
  graceful-degrading fallback may suffice. The operational constraint
  thus directly sets the economic floor.
\end{itemize}

\section{What We Still Don't Know}\label{what-we-still-dont-know}

\begin{itemize}
\tightlist
\item
  \textbf{Tokenizer exactness for technical domains}: While the ``one
  token ≈ 4 characters'' rule of thumb is convenient, enterprise Q\&A
  prompts may contain code, JSON, or domain‑specific terminology that
  tokenizes differently. The exact per‑token count for a given prompt
  distribution is {[}VERIFY{]} and would require running the model's
  tokenizer on a representative sample of real user queries.
\item
  \textbf{KV cache scaling with longer contexts}: The canonical scenario
  uses \textasciitilde9.2 K input tokens. If retrieval quality degrades
  and context lengths grow to 32 K or 128 K, the KV cache memory per
  request grows linearly, and the single‑host FP16 architecture would no
  longer suffice without quantization or model parallelism. The
  breakpoint is {[}VERIFY{]} and workload‑dependent.
\item
  \textbf{Impact of reranking on latency}: The TTFT budget assumes
  retrieval (\textasciitilde120 ms) plus prefill. If a late‑stage
  reranker is added to the pipeline, the retrieval component's latency
  may increase, eating into the TTFT budget and potentially requiring a
  wider SLO margin or a faster retriever. The magnitude of this effect
  is {[}HYPOTHESIS{]} and should be measured before production
  deployment.
\item
  \textbf{Economic durability at scale}: The tokens‑per‑dollar numbers
  are derived from 2026 on‑demand H100 pricing. Spot instances, reserved
  contracts, or custom cloud agreements could shift the economics
  significantly. The durability of the economic model across provider
  changes is {[}HYPOTHESIS{]}.
\end{itemize}

\section{End-of-Chapter Mini-Case}\label{end-of-chapter-mini-case}

An architect is brought into the early design of an internal Q\&A
platform. The stakeholder says: ``We have thousands of employees who
want to ask questions over our internal documents. We need it to be
accurate and fast, but we don't know how many thousands or how fast is
fast enough.'' Before any architecture can be defended, the architect
does the workload characterization that this chapter walks through.

From the token layer alone (as we did in Ch. 4), the architect can
already establish: the unit is tokens; the request shape will be prompt
+ retrieved context + output; the workload is input‑heavy; and the first
number to lock down is tokens‑per‑request, because every downstream
decision (which model fits, how much memory, what latency is possible)
is priced against it. Using the canonical scenario as a starting point
--- \textasciitilde2,000 registered users, \textasciitilde5\%
concurrency, \textasciitilde10 rps, \textasciitilde9.2 K input + 300
output tokens --- the architect projects the enterprise's actual
headcount. If the company has 5,000 employees and expects 10\%
concurrent activity during peak Q\&A periods (after a policy rollout),
the concurrency rises to 500 users. With the same \textasciitilde10 s
request duration, throughput climbs to \textasciitilde50 rps, which
would require \textasciitilde5 hosts of 8×H100 each, at a monthly
infrastructure cost of \textasciitilde\$12,500. The architect can now
speak the system's currency --- tokens, rps, dollars per million --- and
can engage the rest of the team on model selection, serving
configuration, and TCO with concrete numbers rather than vague
assurances. The specific sizing --- turning ``thousands of employees''
into a characterized workload across all six dimensions --- is the
payoff of Ch. 4's framework, and it is the prerequisite for every
architecture question that follows.
